\section{Limitations and Future Work}\label{s:8}

Despite promising reductions in queue length and network-wide accumulated delay, several methodological and experimental limitations must be addressed before full-scale benchmarking or real-world deployment:

\begin{enumerate}
\item \textbf{Short evaluation horizon and single-seed benchmarking:} Although training configurations and \texttt{selected\_scenario.csv} specify $3600\,\mathrm{s}$ episodes, the current evaluation outputs in \texttt{outputs/results/} were collected over $100\,\mathrm{s}$ rollouts (20 control steps) with a single evaluation seed (\texttt{seed=42}, one episode per scenario). As demonstrated in Section~\ref{s:7}, short $100\,\mathrm{s}$ windows induce right-censoring in \texttt{tripinfo.xml} metrics and prevent assessment of long-horizon queue dissipation, spillback recovery, and steady-state saturation. Moreover, single-seed evaluation precludes Student-$t$ confidence interval estimation across seeds, and supplementary \texttt{spillback} and \texttt{recovery\_time} fields remain unpopulated. Future evaluation will execute full $3600\,\mathrm{s}$ episodes across multiple random seeds (e.g., 5 seeds per scenario) to establish statistical significance.

\item \textbf{Incomplete baseline and algorithm coverage:} While the software framework defines \texttt{PPO}, \texttt{ActuatedController}, and \texttt{MaxPressureController}, the actuated and max-pressure baselines remain unimplemented stubs (\texttt{NotImplementedError}), and the current benchmark reports only \texttt{dqn\_agent} versus \texttt{fixed\_time}. Furthermore, the \texttt{fixed\_time} baseline executes the static \texttt{programID="1"} defined in \texttt{roundabout.net.xml} ($31\,\mathrm{s}$ uniform phase durations across 32 bitmask states, including all-red and all-green phases) rather than a Webster- or demand-optimized fixed-time signal plan, and logs a constant dummy action (\texttt{0}) so its internal phase transitions are not reflected in \texttt{phase\_switch\_rate}. Benchmarking DQN and PPO against tuned Webster fixed-time, actuated, and max-pressure controllers is essential for rigorous comparative validation.

\item \textbf{Action-space granularity and unpenalized switching:} Under \texttt{programID="1"}, the agent selects among 32 phase combinations ($2^5$ bitmasks of the five controlled approaches), with \texttt{red\_time=0} (all-red clearance disabled) and \texttt{enforce\_max\_green=false} in \texttt{roundabout.yaml}, alongside $w_{\text{flicker}}=0.0$ in \texttt{dqn\_roundabout.yaml}. The resulting high switch request rate ($0.58$ per step) may increase clearance overhead and emissions, while unpruned phase combinations (such as all-green \texttt{GGGGGGGGGG} or all-red \texttt{rrrrrrrrrr}) enlarge the exploration space. Restricting the action set to conflict-safe green phases, enabling non-zero flicker regularization, and enforcing maximum-green bounds will improve stability and physical realism.

\item \textbf{Lack of explicit emergency-vehicle preemption:} In Scenario~80 (\texttt{80\_emergency\_multi\_ambulance}, 8.0\% emergency vehicles), ambulances are tracked in the state vector and assigned $w_{\text{pcu}}=1.0$ in \texttt{CombinedReward}, but receive no priority weight over passenger cars. Incorporating class-priority delay weights or explicit ambulance preemption constraints into the reward and signal-transition logic is an important extension for safety-critical urban corridors.

\item \textbf{Isolated intersection scope and sim-to-real gap:} The controller optimizes the Ly Thai To roundabout in isolation without modeling upstream metering or downstream arterial spillback to adjacent District~3 and District~10 junctions. In addition, although SUMO's \texttt{SL2015} sublane model captures lateral filtering at $1.2\,\mathrm{m}$ resolution, real-world Vietnamese mixed traffic exhibits informal behaviors (red-light encroachment, unmodeled pedestrian crossings, roadside parking obstructions, and noisy camera/loop detection) absent from idealized TraCI observations. Future work will extend the architecture to multi-agent corridor coordination, domain randomization with noisy observations, and hardware-in-the-loop validation.
\end{enumerate}
